% 褐矮星（综述）
% license CCBYSA3
% type Wiki

本文根据 CC-BY-SA 协议转载翻译自维基百科\href{https://en.wikipedia.org/wiki/Brown_dwarf}{相关文章}。

\begin{figure}[ht]
\centering
\includegraphics[width=6cm]{./figures/6432b0e71e12ceea.png}
\caption{T 型褐矮星的艺术概念图} \label{fig_Brdw_1}
\end{figure}
\begin{figure}[ht]
\centering
\includegraphics[width=6cm]{./figures/dc30c14f9295e490.png}
\caption{对比：大多数褐矮星的体积仅比木星略大（约大15–20\%）$^{[1]}$，但由于其密度更高，质量仍可达到木星的 80 倍。图像按比例绘制，其中木星的半径约为地球的 11 倍，而太阳的半径约为木星的 10 倍。} \label{fig_Brdw_2}
\end{figure}
褐矮星是亚恒星天体，其质量大于最大的气态巨行星，但小于质量最小的主序星。它们的质量大约为木星的 13–80 倍（($M_J$)）$^{[2]}$——不足以在其核心中维持将氢聚变为氦的核反应，但又足够巨大，可以通过氘（$^2\mathrm{H}$）的聚变释放一定的光和热。氘是氢的一种同位素，含有一个中子和一个质子，其聚变所需温度更低。质量最大的褐矮星（$(>65,M_J)$）还能够进行锂（$(^7\mathrm{Li})$）的聚变。

天文学家按照光谱型对自发光天体进行分类，这一分类与表面温度密切相关。褐矮星分布在 M 型（2100–3500 K）、L 型（1300–2100 K）、T 型（600–1300 K）和 Y 型（<600 K） 光谱类别中 $^{[3][4]}$。由于褐矮星不会发生稳定的氢聚变，它们会随着时间推移逐渐冷却，并在演化过程中依次进入更晚的光谱型。

“褐矮星”中的“褐色”原意是指介于红色与黑色之间的一种颜色 $^{[5]}$。以肉眼观察，大多数褐矮星呈现为品红色，而其他褐矮星则会因温度不同而显示出不同的颜色 $^{[3][6]}$。褐矮星可能是完全对流的，内部不存在明显的分层结构或随深度变化的化学分异 $^{[7]}$。

尽管早在 20 世纪 60 年代就有人从理论上预言了褐矮星的存在，但直到 1994 年，第一批明确无误的褐矮星才被发现 $^{[8]}$。由于褐矮星的表面温度相对较低，它们在可见光波段并不明亮，而主要在红外波段辐射能量。随着红外探测设备能力的提升，天文学家已经识别出数以千计的褐矮星。目前已知距离最近的褐矮星位于 Luhman 16 系统，这是一个由 L 型和 T 型褐矮星组成的双星系统，距离太阳约 6.5 光年（2.0 秒差距）。在距离上，Luhman 16 是继半人马座 $\alpha$ 星系和巴纳德星之后，距离太阳第三近的恒星系统。
\subsection{历史}
\begin{figure}[ht]
\centering
\includegraphics[width=6cm]{./figures/16619a30e2a82cab.png}
\caption{较小的天体是葛利泽 229B（Gliese 229B），其质量约为木星的 20 至 50 倍，绕恒星葛利泽 229 运行。它位于天兔座，距离地球约 19 光年。} \label{fig_Brdw_3}
\end{figure}
\subsubsection{早期理论研究}
\begin{figure}[ht]
\centering
\includegraphics[width=6cm]{./figures/a5a243a9e6a937ce.png}
\caption{行星、褐矮星、恒星（非按比例）} \label{fig_Brdw_4}
\end{figure}
20 世纪 60 年代，Shiv Kumar 提出了后来被称为褐矮星的天体的存在理论；它们最初被称为黑矮星$^{[9]}$，这一名称用于指在太空中自由漂浮、质量不足以维持氢聚变的暗弱亚恒星天体。然而：\\

(a) “黑矮星”一词此前已被用来指冷却后的白矮星；\\
(b) 红矮星能够进行氢聚变；\\
(c) 这些天体在其生命早期可能在可见波段仍然具有一定亮度。\\

因此，人们提出了其他名称来称呼这些天体，包括 planetar 和 substar。1975 年，Jill Tarter 在加州大学伯克利分校的博士论文中，首次建议使用“褐矮星（brown dwarf）”这一术语，用“褐色”来表示一种“介于红色与黑色之间”的颜色，意指这些矮星看起来昏暗、黯淡、无光泽，尽管它们并非真正呈褐色$^{[5][10][11]}$。

“黑矮星”这一术语至今仍用于指已经冷却到几乎不再发出可探测光辐射的白矮星。然而，即便是最低质量的白矮星，要冷却到这一温度所需的时间也被计算为超过宇宙当前的年龄，因此这类天体预计至今尚未存在$^{[12]}$。

关于最低质量恒星及氢燃烧下限的早期理论认为：第一星族天体若质量小于 $0.07,M_\odot$，或第二星族天体若质量小于 $0.09,M_\odot$，将永远不会经历正常的恒星演化过程，而会成为完全简并的恒星$^{[13]}$。由此产生的褐矮星有时被称为“失败的恒星”$^{[14]}$。首个自洽的氢燃烧最小质量计算结果证实，第一星族天体的下限约为 $0.07-0.08,M_\odot$,$^{[15][16]}$。
\subsubsection{氘聚变}
20 世纪 80 年代末，人们发现质量低至 $0.013,M_\odot$（约 $13.6,M_J$）的天体也可以发生氘聚变，同时褐矮星冷却外层大气中尘埃形成的影响也被认识到，这些发现对上述理论提出了挑战。然而，这类天体极难被发现，因为它们几乎不在可见光波段发射辐射，其最强的辐射位于红外（IR）波段。当时的地基红外探测器精度不足，难以可靠识别褐矮星。

此后，研究者采用多种方法开展了大量搜索工作，包括：对场星的多色成像巡天；对主序矮星和白矮星暗弱伴星的成像巡天；对年轻星团的调查；以及通过径向速度监测寻找近距离伴星。
\subsubsection{GD 165B 与 L 型}

多年来，发现褐矮星的尝试一直未获成功。直到 1988 年，在一次针对白矮星的红外巡天中，人们发现白矮星 GD 165 周围存在一个暗弱伴星。该伴星 GD 165B 的光谱呈现出非常红且难以解释的特征，既不符合低质量红矮星的典型特征，也不同于当时已知的最冷 M 型矮星。这表明，GD 165B 必须被归类为一种比已知 M 型矮星更冷的天体。

在接近十年的时间里，GD 165B 一直是独一无二的例子。直到 1997 年 2MASS（Two Micron All-Sky Survey）启动后，人们才发现了大量在颜色和光谱特征上与其相似的天体。

如今，GD 165B 被认为是后来被称为 L 型矮星的一类天体的原型$^{[17][18]}$。

尽管当时发现如此冷的矮星具有重大意义，但学界一度争论 GD 165B 究竟应被归类为褐矮星，还是仅仅是一颗极低质量恒星，因为在观测上二者极难区分$^{[19][20]}$。

在 GD 165B 被发现之后不久，又陆续有其他褐矮星候选体被报告。然而，其中大多数未能通过检验，因为它们的大气中缺乏锂元素，表明它们实际上是恒星。真正的恒星会在约$10^8$年左右燃尽其锂，而褐矮星（尽管其温度和光度可能与恒星相近）则不会发生这一过程。因此，在年龄超过 $100,\mathrm{Myr}$ 的天体大气中探测到锂元素，几乎可以确定该天体是一颗褐矮星。
\subsubsection{Teide 1 与 M }
首个被确认的 M 型褐矮星于 1994 年由西班牙天体物理学家 Rafael Rebolo（课题组负责人）、María Rosa Zapatero-Osorio 和 Eduardo L. Martín 发现$^{[21]}$。该天体位于昴宿星团这一疏散星团中，被命名为 Teide 1。相关发现论文于 1995 年 5 月投稿至《自然》，并于 1995 年 9 月 14 日正式发表$^{[22][23]}$。在该期《自然》杂志的封面上，醒目标题写道：“褐矮星被发现，官方确认”。

Teide 1 是 IAC 团队于 1994 年 1 月 6 日在特内里费岛泰德天文台，使用 80 厘米口径望远镜（IAC 80）获取的图像中发现的；其光谱则于 1994 年 12 月首次由拉帕尔马岛罗克·德·洛斯·穆查乔斯天文台的 4.2 米威廉·赫歇尔望远镜观测获得。由于 Teide 1 属于年轻的昴宿星团，其距离、化学组成和年龄得以较为准确地确定。利用当时最先进的恒星与亚恒星演化模型，研究团队估计 Teide 1 的质量为$55 \pm 15,M_J$,$^{[24]}$，低于恒星质量下限。此后，Teide 1 成为年轻褐矮星研究中的一个重要参照对象。

从理论上讲，质量低于$65,M_J$的褐矮星在其整个演化过程中都无法通过热核聚变燃烧锂。这一事实构成了锂检验原理的基础之一，用以判断低光度、低表面温度天体是否具有亚恒星性质。

1995 年 11 月，使用 Keck I 望远镜获得的高质量光谱数据显示，Teide 1 仍然保留着形成昴宿星团恒星的原始分子云中的初始锂丰度，从而证明其核心中不存在热核聚变过程。这些观测结果不仅完全确认了 Teide 1 是一颗褐矮星，也验证了光谱锂检验方法的有效性$^{[25]}$。

在相当一段时间内，Teide 1 是通过直接观测所识别的、太阳系外已知质量最小的天体。此后，人们已鉴认出超过 1,800 颗褐矮星$^{[26]}$，其中一些距离地球非常近，例如 Epsilon Indi Ba 与 Bb——一对与一颗类太阳恒星引力束缚、距离太阳约 12 光年的褐矮星系统$^{[27]}$，以及 Luhman 16——一个距离太阳仅 6.5 光年的褐矮星双星系统。
\subsubsection{Gliese 229B 与 T 型}
首个 T 型褐矮星于 1994 年由加州理工学院的天文学家 Shrinivas Kulkarni、中岛敬史（Tadashi Nakajima）、Keith Matthews 和 Rebecca Oppenheimer，以及约翰斯·霍普金斯大学的科学家 Samuel T. Durrance 和 David Golimowski 共同发现$^{[28]}$。该天体于 1995 年被确认是 Gliese 229 的一颗亚恒星伴星，即 Gliese 229b。Gliese 229b 与 Teide 1 一同构成了最早两例清晰的褐矮星观测证据。二者均在 1995 年通过探测到 670.8 nm 的锂谱线得到确认，其中 Gliese 229b 的温度和光度显著低于恒星范围。

其近红外光谱清楚地显示出 2 微米处的甲烷吸收带，这一特征此前仅在巨行星大气以及土星卫星泰坦的大气中被观测到。甲烷吸收并不可能出现在任何主序恒星的温度条件下。正是这一发现促成了比 L 型矮星更冷的新光谱类型——T 型矮星的确立，而 Gliese 229b 则被视为该类型的原型。
\subsection{理论}
\begin{figure}[ht]
\centering
\includegraphics[width=6cm]{./figures/851acc2df0619164.png}
\caption{} \label{fig_Brdw_5}
\end{figure}
恒星形成的标准机制是：一团寒冷的星际气体与尘埃云在自身引力作用下发生坍缩。随着云团收缩，由于 Kelvin–Helmholtz 机制，其温度会升高。在这一过程的早期阶段，收缩中的气体会迅速将大量能量以辐射形式释放出去，从而使坍缩得以继续。最终，中心区域的密度增大到足以俘获辐射的程度。于是，坍缩云团中心的温度和密度会随时间显著上升，收缩速度逐渐减慢，直到核心条件变得足够高温、高密，从而在原恒星核心中触发热核反应。对于一颗典型恒星而言，其核心内由热核聚变反应产生的气体压和辐射压可以支撑恒星抵抗进一步的引力收缩。此时达到流体静力平衡，恒星将在其大部分寿命中以主序星的形式，将氢聚变为氦。

然而，如果原恒星的初始质量$^{[29]}$小于约 (0.08,M_{\odot})$^{[30]}$，则正常的氢热核聚变反应将无法在其核心中点燃。引力收缩并不能非常有效地加热这样一个低质量的原恒星；在核心温度尚未升高到足以触发聚变之前，密度便已上升到电子彼此紧密堆积的程度，从而产生量子意义上的电子简并压。根据褐矮星内部结构模型，核心中典型的密度、温度和压力条件预计如下：
\begin{itemize}
\item 
\[
10,\mathrm{g/cm^3}\ \lesssim\ \rho_c\ \lesssim\ 10^3,\mathrm{g/cm^3},~
\]
\item 
\[
T_c \lesssim 3\times 10^6,\mathrm{K},~
\]
\item 
\[
P_c \sim 10^5,\mathrm{Mbar}.~
\]
\end{itemize}
这意味着，该原恒星在任何时候都不具备维持氢聚变所需的质量或密度条件。正在落入的物质会被电子简并压所阻止，无法达到所需的高密度和高压力。

进一步的引力收缩因此被抑制，最终形成的结果是一颗褐矮星，它仅通过向外辐射自身内部的热能而逐渐冷却。需要注意的是，从原理上讲，褐矮星有可能在不点燃氢聚变的情况下，缓慢吸积物质并使其质量超过氢燃烧下限。这种情况可能通过褐矮星双星系统中的质量转移过程而发生$^{[29]}$。
\subsubsection{高质量褐矮星与低质量恒星}
锂通常存在于褐矮星中，而在低质量恒星中则一般不存在。恒星一旦达到足以发生氢聚变的高温，其内部的锂会很快被消耗殆尽。锂-7 与质子的聚变会生成两个氦-4 原子核，而该反应所需的温度仅略低于氢聚变所需的温度。低质量恒星中的对流作用会使整个恒星内部的锂最终全部被耗尽。因此，在一个候选褐矮星的光谱中探测到锂谱线，是其确为次恒星天体的一个强有力指标。

\textbf{锂检验}

褐矮星可以分为两类：一类具有足够的质量可以发生锂聚变，另一类则不能。这一区分方法被称为锂检验$^{[31]}$。

像太阳这样质量更大的恒星，也可能在其外层保留锂，因为这些外层从未达到足以发生锂聚变的温度，而且其对流层并不会与核心混合；若发生混合，锂将会迅速被消耗。这类较大质量的恒星可以通过其尺寸和光度轻易地区分于褐矮星。

相反地，处在质量上限附近的褐矮星，在其年轻阶段可能足够炽热，从而耗尽自身的锂。质量大于 $65,M_J$的褐矮星，在其约五亿年寿命之前便可以燃烧掉锂$^{[32]}$。

\textbf{大气中的甲烷}

与恒星不同，年龄较老的褐矮星有时会冷却到足够低的温度，使其大气在极长的时间尺度上积累可观测量的甲烷，而甲烷无法在更高温的天体中形成。通过这种方式得到确认的褐矮星包括 Gliese 229B。

\textbf{铁、硅酸盐与硫化物云}

主序星在演化过程中会逐渐冷却，但最终会达到一个能够通过稳定核聚变维持的最小总光度。该光度因恒星而异，但通常至少为太阳光度的 0.01\%$^{[\text{citation needed}]}$。而褐矮星则在其整个寿命中持续冷却并变暗；足够古老的褐矮星将会暗淡到无法被探测的程度。
\begin{figure}[ht]
\centering
\includegraphics[width=8cm]{./figures/963216b7bed1d372.png}
\caption{早期 T 型褐矮星 SIMP J0136+09 和 2MASS J2139+02（左侧两个面板）的云模型，以及晚期 T 型褐矮星 2M0050–3322 的云模型。} \label{fig_Brdw_6}
\end{figure} 
云被用来解释晚期 L 型褐矮星中铁氢化物（FeH）光谱线的减弱。铁云会使上层大气中的 FeH 被耗尽，同时云层遮挡了对仍然含有 FeH 的更深层大气的观测视线。随后，在中期到晚期 T 型褐矮星更低的温度条件下，这种化学物质再次增强，被解释为云层受到扰动，使望远镜能够观测到仍然含有 FeH 的更深大气层$^{[33]}$。年轻的 L/T 型褐矮星（L2–T4）表现出很强的光变，这可能可以用云层、热点、磁驱动的极光或热化学不稳定性来解释$^{[34]}$。这些褐矮星的云被解释为厚度变化的铁云，或是下方一层厚铁云与上方一层硅酸盐云的组合。上层硅酸盐云可能由石英（quartz）、顽火辉石（enstatite）、刚玉（corundum）和／或橄榄石（forsterite）组成$^{[35][36]}$。然而，目前尚不清楚硅酸盐云是否对所有年轻天体都是必需的$^{[37]}$。硅酸盐吸收可以在 8–12 μm 的中红外波段被直接观测到。利用 Spitzer 的 IRS 进行的观测表明，硅酸盐吸收在 L2–L8 型褐矮星中较为常见，但并非普遍存在$^{[38]}$。此外，MIRI 还在行星质量伴星 VHS 1256b 中观测到了硅酸盐吸收$^{[39]}$。

作为大气对流过程一部分的“铁雨”只可能发生在褐矮星中，而不会出现在小质量恒星中。关于铁雨的光谱研究仍在进行中，但并非所有褐矮星在任何时候都会表现出这种大气异常。2013 年，人们在附近的 Luhman 16 系统中，对其 B 分量周围的非均匀含铁大气进行了成像$^{[40]}$。

对于晚期 T 型褐矮星，仅开展了少量的变光搜索研究。理论预测，在晚期 T 型褐矮星中会形成由铬、氯化钾以及多种硫化物构成的薄云层，这些硫化物包括硫化锰、硫化钠和硫化锌$^{[41]}$。对可变的 T7 型褐矮星 2M0050–3322 的解释认为，其最上层为氯化钾云，中间层为硫化钠云，下层为硫化锰云。上两层云的斑块状分布可以解释为何甲烷和水汽吸收带呈现出变化$^{[42]}$。

在最低温度的 Y 型褐矮星 WISE 0855−0714 中，硫化物云和水冰云的斑块状云层可能覆盖其约 50\% 的表面$^{[43]}$。
\subsubsection{低质量褐矮星与高质量行星的比较}
\begin{figure}[ht]
\centering
\includegraphics[width=6cm]{./figures/e55ea39cbe25dd2e.png}
\caption{} \label{fig_Brdw_7}
\end{figure}
与恒星类似，褐矮星也是独立形成的；但与恒星不同的是，它们缺乏足够的质量来“点燃”氢聚变。和所有恒星一样，褐矮星既可以单独存在，也可以与其他恒星近距离共存。有些褐矮星环绕恒星运行，并且像行星一样，其轨道也可能具有较大的偏心率。

\textbf{尺寸与燃料燃烧的不确定性}

褐矮星的半径大致都与木星相当。在其质量范围的高端（60–90 $M_{\mathrm{J}}$），褐矮星的体积主要受电子简并压控制，这一点与白矮星类似$^{[44]}$；而在质量范围的低端（约 10 $M_{\mathrm{J}}$），其体积主要受库仑压控制，这又与行星相似。综合来看，在整个可能的质量范围内，褐矮星的半径变化仅有约 10–15\%。此外，质量–半径关系在从约一个土星质量到氢燃烧起始点（$0.080\pm0.008,M_\odot$）之间几乎没有变化，这表明从这一角度看，褐矮星只是高质量的类木行星$^{[45]}$。这也使得在观测上区分褐矮星与行星变得十分困难。

另外，许多褐矮星根本不发生任何核聚变；即便是处在质量上限（超过 60 $M_{\mathrm{J}}$）的褐矮星，也会因冷却速度很快而在约一千万年后不再进行聚变反应。

\textbf{热辐射光谱}

X 射线和红外光谱是识别褐矮星的重要特征。一些褐矮星会发射 X 射线；而所有“温暖”的褐矮星都会在红光和红外波段持续发光，直到它们冷却到类似行星的温度（低于 1000 K）。

气态巨行星也具有某些与褐矮星相似的特性。与太阳类似，木星和土星主要由氢和氦组成。尽管土星的质量只有木星的约 30\%，其体积却几乎与木星相当。太阳系中的四颗巨行星中，有三颗（木星、土星和海王星）辐射出的热量远多于它们从太阳接收到的热量，最高可达约两倍$^{[46][47]}$。这四颗巨行星都拥有各自的“行星系统”，即由大量卫星构成的广泛卫星体系。

\textbf{现行 IAU 标准}

目前，国际天文学联合会（IAU）将质量高于 13 个木星质量（$13,M_{\mathrm{J}}$，即氘发生热核聚变的质量下限）的天体视为褐矮星，而将质量低于该数值（并且绕恒星或恒星遗迹运行）的天体视为行星。能够触发并维持氢燃烧所需的最小质量（约 $80,M_{\mathrm{J}}$）构成了这一分类的上限$^{[48]}$。

同时，学界也在讨论：与其依据基于核聚变反应的理论质量界限，不如依据形成过程来定义褐矮星是否更为合理$^{[3]}$。在这种解释下，褐矮星是恒星形成过程所产生的最低质量产物，而行星则是在恒星周围的吸积盘中形成的天体。已发现的一些最冷的自由漂浮天体（如 WISE 0855），以及已知质量最低的年轻天体（如 PSO J318.5−22），被认为其质量低于 $13,M_{\mathrm{J}}$，因此由于其究竟应被视为流浪行星还是褐矮星存在歧义，常被称为行星质量天体。此外，也已知存在绕褐矮星运行的行星质量天体，例如 2M1207b、2MASS J044144b 和 Oph 98 B。

将 $13,M_{\mathrm{J}}$ 作为分界线更多是一种经验法则，而非具有精确定义的物理常数。质量更大的天体会燃烧其大部分氘，而质量更小的天体只会燃烧极少量，$13,M_{\mathrm{J}}$ 大致位于两者之间$^{[49]}$。氘燃烧的程度还在一定程度上取决于天体的组成，尤其是氦和氘的含量，以及重元素所占比例；这些因素会影响大气的不透明度，从而影响辐射冷却速率$^{[50]}$。

截至 2011 年，《系外行星百科全书》（Extrasolar Planets Encyclopaedia）收录了质量高达 $25,M_{\mathrm{J}}$ 的天体，并指出：“在观测到的质量谱中，$13,M_{\mathrm{J}}$ 附近并不存在任何特殊特征，这一事实强化了放弃该质量界限的选择”$^{[51]}$。到 2016 年，基于质量–密度关系的研究，该上限被提高至 $60,M_{\mathrm{J}}$,$^{[52][53]}$。

Exoplanet Data Explorer 收录了质量最高达 $24,M_{\mathrm{J}}$ 的天体，并附有说明：“IAU 工作组提出的 $13,M_{\mathrm{J}}$ 区分标准在具有岩石核心的行星情形下缺乏物理动机，并且由于 $\sin i$ 不确定性，在观测上也存在问题”$^{[54]}$。NASA 系外行星档案（NASA Exoplanet Archive）则收录了质量（或最小质量）不超过 $30,M_{\mathrm{J}}$ 的天体$^{[55]}$。

\textbf{次褐矮星}
\begin{figure}[ht]
\centering
\includegraphics[width=6cm]{./figures/8c9a960c9abbf3d5.png}
\caption{太阳、一个年轻的次褐矮星与木星的体积大小对比。随着次褐矮星逐渐老化，它会不断冷却并收缩。} \label{fig_Brdw_8}
\end{figure}
质量低于 13 MJ 的天体被称为次褐矮星（sub-brown dwarfs）或行星质量褐矮星。它们的形成方式与恒星和褐矮星相同（即通过气体云的引力坍缩），但其质量低于氘发生热核聚变的下限 $^{[56]}$。有些研究者将它们称为自由漂浮行星$^{[57]}$，而另一些则称其为行星质量褐矮星 $^{[58]}$。
\subsubsection{其他物理性质在质量估计中的作用}
尽管光谱特征可以帮助区分低质量恒星与褐矮星，但在许多情况下仍需要通过质量估计才能得出结论。质量估计背后的理论认为，质量相近的褐矮星具有相似的形成过程，并且在形成初期温度较高。有些天体的光谱型与低质量恒星相似，例如 2M1101AB。随着时间推移而冷却，这些褐矮星会根据其质量保留不同的光度范围 $^{[59]}$。

如果缺乏年龄和光度信息，质量估计将变得十分困难。例如，一个 L 型褐矮星可能是质量较高、年龄较大的褐矮星（甚至可能是低质量恒星），也可能是质量极低、但非常年轻的褐矮星。对于 Y 型褐矮星而言，这个问题相对较小，因为即使在较大的年龄估计下，它们仍然是接近次褐矮星质量下限的低质量天体 $^{[60]}$。而对于 L 型和 T 型褐矮星，获得准确的年龄估计仍然十分重要。相比之下，光度并不是最棘手的问题，因为可以通过光谱能量分布来估计 $^{[61]}$。

年龄估计通常有两种方法：其一，褐矮星本身非常年轻，并且仍然保留与年轻阶段相关的光谱特征；其二，褐矮星与某颗恒星或恒星群（如疏散星团或恒星协会）具有共同运动特性，在这些系统中年龄更容易确定。

利用第二种方法研究的一个非常年轻的褐矮星实例是 2M1207 及其伴星 2M1207b。根据其空间位置、自行和光谱特征，该系统被确定属于约 800 万年历史的 TW Hydrae 恒星协会，其次级天体的质量被估计为$8 \pm 2M_J$，低于氘燃烧下限 $^{[62]}$。

而通过共同运动方法得到极高年龄估计的一个例子是褐矮星–白矮星双星系统 COCONUTS-1，其中白矮星的年龄估计为$7.3(^{+2.8}_{-1.6})$十亿年。在这一案例中，质量并非直接由年龄推导，而是通过共同运动获得了精确的距离估计（利用 Gaia 的视差测量）。在此基础上，研究者估计了该褐矮星的半径，并进一步推算其质量为$15.4(^{+0.9}_{-0.8}M_J)$.$^{[63]}$。
\subsection{观测}
\subsubsection{褐矮星的分类}
\textbf{光谱型 M}
\begin{figure}[ht]
\centering
\includegraphics[width=6cm]{./figures/0a732b5953b02bea.png}
\caption{晚期$M$型矮星的艺术想象图} \label{fig_Brdw_9}
\end{figure}
这些是光谱型为 M5.5 或更晚的棕矮星，也被称为晚期$M$型矮星。所有光谱型为 M 的棕矮星都是年轻天体，例如 Teide 1 ——这是发现的第一颗$M$型棕矮星——以及 LP 944-20，它是距离最近的$M$型棕矮星。

\textbf{光谱级 L}
\begin{figure}[ht]
\centering
\includegraphics[width=6cm]{./figures/a812750df0497cb6.png}
\caption{$L$型矮星的艺术概念图} \label{fig_Brdw_10}
\end{figure}
M 型光谱级（经典恒星序列中最冷的一类）的决定性特征，是其光学光谱由二氧化钛（TiO）和二氧化钒（VO）分子的吸收带所主导。然而，白矮星 GD 165 的低温伴星 GD 165B 却完全不具备 M 型矮星典型的 TiO 特征。随后，人们又发现了许多与 GD 165B 类似的天体，这最终促成了一个新的光谱类型——L 型矮星的确立。

L 型矮星在红光光学波段的定义特征，不再是金属氧化物（TiO、VO）的吸收带，而是金属氢化物（FeH、CrH、MgH、CaH）的发射带，以及碱金属原子（Na、K、Rb、Cs）显著的原子谱线。截至 2013 年，已确认的 L 型矮星超过 900 个$^{[26]}$，其中大多数是通过大视场巡天发现的，包括二维微米全天巡天（2MASS）、南天深近红外巡天（DENIS）以及斯隆数字化巡天（SDSS）。

这一光谱类型还包含了最冷的主序恒星（质量大于约 80 MJ），它们的光谱型介于 L2 到 L6 之间$^{[64]}$。

\textbf{光谱级 T}
\begin{figure}[ht]
\centering
\includegraphics[width=6cm]{./figures/b57b3066df5ac4fb.png}
\caption{$T$型矮星的艺术构想图} \label{fig_Brdw_11}
\end{figure}
由于 GD 165B 是 L 型矮星的原型，Gliese 229B 则成为第二个新光谱类型—— T 型矮星的原型。T 型矮星呈粉紫色（洋红色）外观。与 L 型矮星在近红外（NIR）光谱中表现出强烈的水$H(_2)O$和一氧化碳（CO）吸收带不同，Gliese 229B 的近红外光谱主要由甲烷$CH(_4)$吸收带主导；这种特征在太阳系中仅存在于气态巨行星及土卫六（Titan）的大气中。$CH(_4)$、$H(_2)O$ 以及分子氢（$H(_2)$）的碰撞诱导吸收（CIA）共同赋予 Gliese 229B 蓝色的近红外颜色。其红光学波段的光谱陡峭倾斜，并且缺乏$L$型矮星所特有的 FeH 和 CrH 吸收带，而是受到碱金属钠（Na）和钾（K）异常宽广吸收特征的显著影响。

正是由于这些差异，J. Davy Kirkpatrick 提出了$T$光谱型，用于描述在$H$波段和$K$波段中表现出甲烷$CH(_4)$吸收的天体。截至 2013 年，已知的$T$型矮星共有 355 个$^{[26]}$。近年来，Adam Burgasser 和 Tom Geballe 建立了$T$型矮星的近红外光谱分类体系。理论研究表明，$L$型矮星包含极低质量恒星与次恒星天体（褐矮星）的混合体，而$T$型矮星则完全由褐矮星组成。由于$T$型矮星在光谱绿色区域中对钠和钾的强烈吸收，其在人眼视觉下的实际颜色预计并非棕色，而是洋红色$^{[65][66]}$。

早期观测限制了$T$型矮星可被探测的最大距离。然而，如 WISE 0316+4307 这样的$T$型褐矮星，已在距离太阳超过 100 光年的位置被发现。利用 詹姆斯·韦布空间望远镜（JWST）的观测，人们甚至探测到了距离太阳高达 4500 秒差距的$T$型矮星，例如 UNCOVER-BD-1。

\textbf{光谱类型 Y}
\begin{figure}[ht]
\centering
\includegraphics[width=6cm]{./figures/02ecc07011945347.png}
\caption{$Y$型矮星的艺术想象图} \label{fig_Brdw_12}
\end{figure}
2009 年时，已知温度最低的棕矮星其有效温度估计在 500–600 K（227–327 °C；440–620 °F）之间，并被归入光谱型 T9。其代表性例子包括棕矮星 CFBDS J005910.90–011401.3、ULAS J133553.45+113005.2 以及 ULAS J003402.77−005206.7$^{[67]}$。这些天体的光谱在约 1.55 微米处显示出吸收峰$^{[67]}$。Delorme 等人提出，该特征源于氨的吸收，并认为这可作为$T$型向$Y$型过渡的标志，从而将这些天体归类为 Y0 型$^{[67][68]}$。然而，这一特征很难与水和甲烷的吸收区分开来$^{[67]}$，也有其他作者指出，将其直接定为 Y0 型仍为时过早$^{[69]}$。
\begin{figure}[ht]
\centering
\includegraphics[width=6cm]{./figures/0bd389856fb9037f.png}
\caption{WISE 0458+6434 是由 WISE 发现的第一颗超冷棕矮星（绿色圆点）。其中呈现的绿色和蓝色来自红外波段，这些红外波长被映射为可见颜色。} \label{fig_Brdw_13}
\end{figure}
\textbf{提出的光谱类型 H}

2025 年，天文学家 Kevin Luhman 和 Catarina Alves de Oliveira 提出了一种新的光谱类型$H$（$H$源自 hydrocarbon，即“烃类”）。他们利用詹姆斯·韦布空间望远镜的数据，在恒星形成区 IC 348 中识别出许多棕矮星，这些天体的质量非常低（其中许多低于氘燃烧极限），并且在 3.4 μm 处呈现出一条吸收线，该吸收线对应于一种尚未被明确识别的脂肪族烃。这种吸收线将作为$H$型光谱的定义特征。$^{[71]}$

\textbf{垂直混合的作用}
\begin{figure}[ht]
\centering
\includegraphics[width=6cm]{./figures/2efdd4bc98a101d2.png}
\caption{连接一氧化碳与甲烷的主要化学反应路径。短寿命自由基以圆点标注。改编自 Zahnle 与 Marley。$^{[72]}$} \label{fig_Brdw_14}
\end{figure}
在以氢为主的大气中，棕矮星内部存在着一氧化碳（CO）与甲烷（CH₄）之间的化学平衡。在正向反应中，一氧化碳与氢分子反应，生成甲烷和羟基自由基（OH）。随后，羟基自由基可能再与氢反应，形成水分子。在反向反应中，甲烷与羟基反应，重新生成一氧化碳和氢。

在较高温度和较低压力条件下（对应$L$型矮星），该化学反应更倾向于一氧化碳；而在较低温度和较高压力条件下（对应$T$型矮星），反应则更倾向于甲烷，在$T/Y$边界附近由甲烷占主导地位。然而，大气中的垂直混合会改变这种平衡：甲烷可能下沉到更深层的大气，而一氧化碳则会从更深、更高温的层次上升到上部大气。由于 C–O 键的断裂存在能垒，一氧化碳重新转化为甲烷的反应速度较慢，这使得棕矮星的可观测大气处于一种化学非平衡状态。

因此，$L/T$转变主要被定义为：从$L$型矮星中以一氧化碳为主的大气，过渡到$T$型矮星中以甲烷为主的大气。垂直混合的强弱会将这一转变推向更低或更高的温度。这一点对于表面重力较小、且具有扩展大气层的天体尤为重要，例如巨型系外行星。在这些行星中，$L/T$ 转变会发生在更低的温度下；而对于棕矮星而言，这一转变通常发生在约 1200 K 左右。相比之下，系外行星 HR 8799c 的温度约为 1100 K，却并未显示出甲烷特征。$^{[72]}$

$T$型与$Y$型矮星之间的转变通常被定义在约 500 K 左右，这主要是由于这些寒冷且极其昏暗的天体缺乏足够的光谱观测数据。未来借助詹姆斯·韦布空间望远镜（JWST）以及超大望远镜（ELTs）的观测，有望显著扩展具有光谱观测的$Y$型矮星样本。$Y$型矮星的大气通常由强烈的甲烷和水汽吸收特征所主导，并可能还存在氨以及水冰的吸收特征。$^{[73]}$

此外，垂直混合、云层结构、金属丰度、光化学过程、闪电、撞击激波以及金属催化作用等因素，都可能影响$L/T$转变和$T/Y$转变所发生的具体温度。$^{[72]}$

\textbf{次级特征}

年轻的棕矮星具有较低的表面重力，这是因为与具有相似光谱型的场星相比，它们半径更大而质量更低。这类天体通常用后缀字母来标注其表面重力特征：$\beta$（beta）表示中等表面重力，γ（gamma）表示低表面重力。低表面重力的判据包括：CaH、K I 和 Na I 吸收线较弱，以及VO 吸收线较强。$^{[76]}$α（alpha）表示正常表面重力，通常会被省略不写。有时，极低的表面重力会用δ（delta）来表示。$^{[78]}$后缀 “pec”表示 peculiar（异常），该后缀仍用于描述其他不寻常的特征，通常是对多种性质的综合概括，可指示低表面重力、次矮星（subdwarf）或未分辨的双星系统等情况。$^{[79]}$前缀 “sd”表示次矮星，仅用于冷次矮星，该前缀意味着低金属丰度，并且其运动学性质更接近银河晕星而非银盘星。$^{[75]}$ 次矮星在颜色上通常比银盘天体更偏蓝。$^{[80]}$红色后缀用于描述颜色偏红的天体，这通常与较老的年龄有关，而不被解释为低表面重力，而是与较高的尘埃含量有关。$^{[77]}$ $^{[78]}$蓝色后缀则描述近红外颜色偏蓝、且不能用低金属丰度来解释的天体。其中一部分被解释为$L+T$型双星系统，而另一些并非双星，例如 2MASS J11263991−5003550，其性质被认为与稀薄和/或大颗粒云层有关。$^{[78]}$
\subsubsection{棕矮星的光谱与大气性质}
\begin{figure}[ht]
\centering
\includegraphics[width=6cm]{./figures/ebca80dcfaaa1db9.png}
\caption{棕矮星内部结构的艺术示意图。由于层状位移，特定深度处的云层出现错位。} \label{fig_Brdw_15}
\end{figure}
$L$型和$T$型矮星所辐射的绝大部分能量通量集中在 1–2.5 微米的近红外波段。沿着晚期$M$型、$L$型到$T$型矮星序列，温度不断降低并持续下降，这使得近红外光谱极为丰富，包含多种多样的谱线特征：既有相对狭窄的中性原子谱线，也有宽广的分子吸收带，而这些特征对温度、表面重力以及金属丰度的依赖关系各不相同。此外，这种低温环境有利于气体发生凝结并形成固态颗粒。
\begin{figure}[ht]
\centering
\includegraphics[width=6cm]{./figures/6c041779a34cc547.png}
\caption{测量到的风（斯皮策空间望远镜；艺术概念图；2020 年 4 月 9 日）$^{[81]}$} \label{fig_Brdw_16}
\end{figure}
已知褐矮星的典型大气温度范围约为 2200 K 至 750 K$^{[65]}$。与依靠稳定内部核聚变自我加热的恒星不同，褐矮星会随时间迅速冷却；质量较大的褐矮星冷却得比质量较小的更慢。有一些证据表明，在光谱型$L$与$T$的过渡区（约 1000 K）附近，褐矮星的冷却过程会减缓$^{[82]}$。

对已知褐矮星候选体的观测揭示了一种红外辐射周期性变亮与变暗的模式，这表明相对较冷且不透明的云层结构遮蔽着在强烈风场搅动下的炽热内部。这类天体上的“天气”被认为极其猛烈，其强度可与、甚至远超木星著名的风暴。

2013 年 1 月 8 日，天文学家利用 NASA 的哈勃空间望远镜和斯皮策空间望远镜，对一颗名为 2MASS J22282889–4310262 的褐矮星的多风暴大气进行了探测，绘制出了迄今为止最为详细的褐矮星“天气图”。该图显示了由风驱动的、行星尺度的云层。上述研究为更深入理解褐矮星以及太阳系外行星的大气性质迈出了重要一步$^{[83]}$。

2020 年 4 月，科学家报告称在附近的一颗褐矮星 2MASS J10475385+2124234 上测得风速为 $+650 \pm 310$ m/s（最高可达约 1450 英里/小时）。为获得这一结果，研究人员将由亮度变化推断出的气象特征旋转运动，与褐矮星内部产生的电磁旋转进行对比。该结果验证了此前关于褐矮星具有高速风场的预测。科学家们期望，这种对比方法可用于研究其他褐矮星及系外行星的大气动力学$^{[84]}$。
\subsubsection{观测技术}
\begin{figure}[ht]
\centering
\includegraphics[width=6cm]{./figures/559b81ff99676a50.png}
\caption{褐矮星 Teide 1、Gliese 229B和WISE 1828+2650与红矮星 Gliese 229A、木星以及太阳的对比} \label{fig_Brdw_17}
\end{figure}
近来，日冕仪（coronagraph）已被用于探测绕明亮可见恒星运行的暗弱天体，其中包括 Gliese 229B。

配备电荷耦合器件（CCD）的高灵敏度望远镜被用于搜索遥远星团中的暗弱天体，例如 Teide 1。

大视场巡天已经识别出一些单个的暗弱天体，如 Kelu-1（距离地球约 30 光年）。

褐矮星常常是在系外行星巡天中被发现的。用于探测系外行星的方法同样适用于褐矮星，而且由于褐矮星本身更亮，它们实际上更容易被探测到。

由于具有强磁场，褐矮星可以成为强烈的射电辐射源。阿雷西博天文台和甚大阵列（VLA）的观测计划已经探测到十多个这样的天体，它们也被称为超冷矮星，因为它们与该类中其他天体共享相似的磁性质$^{[85]}$。对褐矮星射电辐射的探测使得其磁场强度可以被直接测量。
\subsubsection{里程碑}
\begin{itemize}
\item 1995 年：首个褐矮星得到确认。位于昴宿星团中的 Teide 1（一颗 M8 型天体）在西班牙加那利群岛天体物理研究所罗克·德洛斯·穆查乔斯天文台，借助 CCD 被识别出来。
首个甲烷褐矮星得到确认。Gliese 229B 被发现绕红矮星 Gliese 229A（距离约 20 光年）运行；该发现使用了自适应光学日冕仪，以增强加州南部帕洛玛山帕洛玛天文台 60 英寸（1.5 m）反射望远镜的成像效果；随后用 200 英寸（5.1 m）的 Hale 望远镜进行的近红外光谱观测显示其大气中富含甲烷。
\item 1998 年：发现首个发射 X 射线的褐矮星。位于蝘蜓座 I 暗云中的 Cha Helpha 1（M8 型天体）被确定为 X 射线源，其性质与具有对流外层的晚型恒星相似。
\item 1999 年 12 月 15 日：首次在褐矮星中探测到 X 射线耀斑。加州大学的一支团队利用 钱德拉 X 射线天文台监测 LP 944-20（质量约 60 MJ，距离约 16 光年），捕捉到一次持续约 2 小时的耀斑$^{[86]}$。
\item 2000 年 7 月 27 日：首次从褐矮星中探测到射电辐射（包括耀斑态与静态辐射）。甚大阵列（VLA）的一组学生团队探测到 LP 944–20 的射电辐射$^{[87]}$。
\item 2004 年 4 月 30 日：首次在褐矮星周围探测到系外行星候选体：使用 VLT 发现 2M1207b，同时这也是首颗被直接成像的系外行星$^{[88]}$。
\item 2013 年 3 月 20 日：发现距离最近的褐矮星系统：Luhman 16$^{[89]}$。
\item 2014 年 4 月 25 日：发现当时已知最冷的褐矮星。WISE 0855−0714距离地球约 7.2 光年（为距离太阳第七近的恒星系统），其温度介于 −48 ℃ 至 −13 ℃ 之间$^{[90]}$。
\end{itemize}
\subsubsection{褐矮星的 X 射线源}
\begin{figure}[ht]
\centering
\includegraphics[width=8cm]{./figures/0dcede50dee6ee4c.png}
\caption{LP 944-20 在耀斑发生前与耀斑期间的钱德拉（Chandra）X 射线图像} \label{fig_Brdw_18}
\end{figure}

\begin{figure}[ht]
\centering
\includegraphics[width=6cm]{./figures/cb4bae58f97ccb85.png}
\caption{} \label{fig_Brdw_19}
\end{figure}